\pagestyle{empty}
\tight
\begin{tblr}{width=\textwidth,colspec={|Q[c,m,1.9cm]|X[l,m]|},hlines,vlines,hspan=minimal,row{1}={bg=headblue},rowsep=1pt,colsep=3pt}
\SetCell{c,m}\hfil\logo\hfil & \textbf{POLITEKNIK NEGERI MALANG}\par\textbf{JURUSAN TEKNOLOGI INFORMASI}\par\textbf{PROGRAM STUDI: D4 TEKNIK INFORMATIKA} \\
\SetCell[c=2]{c} \textbf{RENCANA PEMBELAJARAN SEMESTER (RPS)} & \\
\end{tblr}
\begin{longtblr}{width=\textwidth,colspec={|Q[l,m,1.9cm]|X[0.85,l,m]|X[1.45,l,m]|X[0.75,l,m]|X[1.15,l,m]|},hlines,vlines,hspan=minimal,rowsep=1pt,colsep=2pt,row{1,3}={bg=headgray,font=\bfseries}}
MATA KULIAH & KODE & BOBOT (SKS)/jam & SEMESTER & TGL. PENYUSUNAN \\
Proyek Sistem Informasi & RTI254009 & 3 SKS/6 jam & 4 & 1 Juli 2025 \\
OTORISASI & Dosen Pengembang RPS & Koordinator RMK & \SetCell[c=2]{l} Ka PRODI &  \\
 & Agung Nugroho Pramudhita, S.T., M.T. &  & \SetCell[c=2]{l}{\vspace{8mm}\par Dr. Ely Setyo Astuti, S.T., M.T.} &  \\
\SetCell[r=19]{l} Capaian Pembelajaran (CP) & \SetCell[c=4]{l,bg=headgray,font=\bfseries} Capaian Pembelajaran Lulusan Program Studi (CPL-Prodi) &  &  &  \\
 & CPL03 & \SetCell[c=3]{l} Mampu mengelola tim, manajemen diri, berkomunikasi secara lisan maupun tertulis dengan baik dan mampu melakukan presentasi. &  &  \\
 & CPL06 & \SetCell[c=3]{l} Mampu merancang, mengimplementasikan, menguji, melakukan deployment, memelihara, serta menjamin mutu perangkat lunak yang menjawab permasalahan dan memenuhi kebutuhan pengguna. &  &  \\
 & CPL08 & \SetCell[c=3]{l} Mampu mengelola proyek teknologi informasi secara profesional melalui perencanaan, koordinasi, komunikasi, dan optimalisasi sumber daya. &  &  \\
 & CPL10 & \SetCell[c=3]{l} Memiliki kompetensi untuk menganalisis persoalan kompleks untuk mengidentifikasi solusi bidang informatika/ilmu komputer dengan mempertimbangkan wawasan ilmu multidisiplin. &  &  \\
 & \SetCell[c=4]{l,bg=headgray,font=\bfseries} Capaian Pembelajaran mata kuliah (CPL-MK) &  &  &  \\
 & CPMK0801 & \SetCell[c=3]{l} Mahasiswa mampu merencanakan proyek sistem informasi dengan estimasi sumber daya, jadwal, dan manajemen risiko. &  &  \\
 & CPMK0802 & \SetCell[c=3]{l} Mahasiswa mampu mengimplementasikan siklus hidup rekayasa perangkat lunak secara menyeluruh dalam pengembangan sistem informasi. &  &  \\
 & CPMK0803 & \SetCell[c=3]{l} Mahasiswa mampu mengeksekusi, memonitor, dan mengendalikan proyek serta mengelola tim secara profesional. &  &  \\
 & CPMK0607 & \SetCell[c=3]{l} Mahasiswa mampu menghasilkan solusi sistem informasi terintegrasi yang menjawab kebutuhan pengguna nyata. &  &  \\
 & CPMK0302 & \SetCell[c=3]{l} Mampu mengelola tim dan proyek secara terencana, mulai dari perencanaan, koordinasi, komunikasi antar anggota, hingga penyampaian hasil kepada pemangku kepentingan. &  &  \\
 & CPMK1008 & \SetCell[c=3]{l} Mampu mengintegrasikan hasil analisis dan desain sistem ke dalam proses pengembangan perangkat lunak secara koheren. &  &  \\
 & \SetCell[c=4]{l,bg=headgray,font=\bfseries} Kemampuan akhir tiap tahapan belajar (Sub-CPMK) &  &  &  \\
 & SCPMK0801-03601 & \SetCell[c=3]{l} Mahasiswa mampu menyusun rencana proyek SI meliputi WBS, jadwal, anggaran, dan risk register. &  &  \\
 & SCPMK0802-03602 & \SetCell[c=3]{l} Mahasiswa mampu mengimplementasikan SI berbasis web menggunakan metodologi SDLC dengan dokumentasi lengkap. &  &  \\
 & SCPMK0803-03603 & \SetCell[c=3]{l} Mahasiswa mampu melaksanakan sprint review, monitoring kemajuan, dan mengelola dinamika tim proyek SI. &  &  \\
 & SCPMK0607-03604 & \SetCell[c=3]{l} Mahasiswa mampu mengintegrasikan front-end, back-end, dan database dalam sistem informasi yang deployable. &  &  \\
 & SCPMK1008-03605 & \SetCell[c=3]{l} Mahasiswa mampu mengintegrasikan hasil analisis kebutuhan dan arsitektur SI ke dalam tech stack dan rencana pengembangan yang koheren. &  &  \\
 & SCPMK0302-03606 & \SetCell[c=3]{l} Mahasiswa mampu mengelola dinamika tim proyek dan menyampaikan hasil kerja tim kepada pemangku kepentingan secara profesional. &  &  \\
\SetCell[r=2]{l} Penilaian & \SetCell[c=4]{l} *Nilai CPMK didapat dari agregasi nilai SUB-CPMK yang dibebankan &  &  &  \\
 & \SetCell[c=4]{l}{\tiny\begin{tblr}{width=\linewidth,colspec={|Q[c,m,0.1\linewidth]|Q[c,m,0.16\linewidth]|X[l,m]|X[l,m]|X[l,m]|X[l,m]|X[l,m]|X[l,m]|Q[c,m,0.08\linewidth]|},hlines,vlines,hspan=minimal,rowsep=1pt,colsep=0.7pt,row{1,2}={bg=headgray,font=\bfseries}}
Kode CPMK & Kode SCPMK & \SetCell[c=6]{c} Bobot per Bentuk Penilaian & & & & & & Total Bobot \\
 & & Dok.Perencanaan & Sprint 1-2 & UTS Milestone & Sprint 3 & Laporan & UAS Demo & \\
CPMK0801 & SCPMK0801-03601 & 15\%: WBS, jadwal, anggaran, risiko & 0\% & 0\% & 0\% & 0\% & 0\% & 15\% \\
CPMK0802 & SCPMK0802-03602 & 0\% & 7\%: implementasi SDLC Sprint 1-2 & 0\% & 10\%: implementasi SDLC Sprint 3 & 0\% & 0\% & 17\% \\
CPMK1008 & SCPMK1008-03605 & 0\% & 3\%: integrasi analisis dan arsitektur ke tech stack & 0\% & 0\% & 0\% & 0\% & 3\% \\
CPMK0803 & SCPMK0803-03603 & 0\% & 0\% & 15\%: monitoring dan review & 0\% & 4\%: laporan kemajuan & 8\%: defense proyek & 27\% \\
CPMK0302 & SCPMK0302-03606 & 0\% & 0\% & 0\% & 0\% & 1\%: koordinasi dan komunikasi tim & 2\%: penyampaian hasil ke stakeholder & 3\% \\
CPMK0607 & SCPMK0607-03604 & 0\% & 10\%: integrasi sistem Sprint 1-2 & 0\% & 10\%: integrasi sistem Sprint 3 & 5\%: laporan akhir & 10\%: demo produk & 35\% \\
\SetCell[c=8]{r}\textbf{Total Per Penilaian} & & & & & & & & \textbf{100\%} \\
\end{tblr}} &  &  &  \\
Diskripsi Singkat Mata Kuliah & \SetCell[c=4]{l} Proyek Sistem Informasi adalah mata kuliah berbasis proyek yang mengintegrasikan kemampuan analisis, desain, implementasi, dan manajemen proyek dalam pengembangan SI berbasis web secara tim menggunakan metodologi agile. &  &  &  \\
Bahan Kajian / Materi Pembelajaran & \SetCell[c=4]{l}\begin{enumerate}\item Perencanaan proyek: WBS, jadwal, anggaran, risiko\item Metodologi SDLC dan Scrum/Agile\item Implementasi SI berbasis web: back-end, front-end, dan database\item Sprint execution, review, dan retrospective\item Deployment, UAT, laporan akhir, dan demo proyek\end{enumerate} &  &  &  \\
\SetCell[r=2]{l} Pustaka & \SetCell[c=4]{l} \textbf{Utama:}\begin{enumerate}\item Sommerville, I. 2016. Software Engineering (10th ed.). Pearson.\item Schwaber, K. 2020. The Scrum Guide. Scrum.org.\item PMI. 2021. PMBOK Guide (7th ed.).\end{enumerate} &  &  &  \\
 & \SetCell[c=4]{l} \textbf{Pendukung:} Materi LMS, dokumentasi resmi framework web, dan jobsheet praktikum. &  &  &  \\
Media Pembelajaran & \SetCell[c=2]{l} \textbf{Software:} IDE, framework web, sistem manajemen proyek, LMS. &  & \SetCell[c=2]{l} \textbf{Hardware:} Komputer/laptop, proyektor. &  \\
Nama Dosen Pengampu & \SetCell[c=4]{l} Tim Pengajar Program Studi D4 TI &  &  &  \\
Matakuliah Syarat & \SetCell[c=4]{l} - &  &  &  \\
\end{longtblr}
\begin{longtblr}{width=\textwidth,colspec={|Q[c,m,0.06\textwidth]|X[1.35,l,m]|X[1.08,l,m]|X[1.16,l,m]|X[0.7,l,m]|X[1.24,l,m]|X[1.06,l,m]|X[0.98,l,m]|Q[c,m,0.05\textwidth]|},hlines,vlines,hspan=minimal,rowhead=2,row{1,2}={bg=headgreen,font=\bfseries},rowsep=1pt,colsep=1.5pt}
Mgg.\par Ke & Kemampuan Akhir Yang Direncanakan (Sub-CP-MK) & Bahan kajian (Materi Pembelajaran) & Bentuk dan Metode Pembelajaran & Estimasi Waktu & Pengalaman Belajar Mahasiswa & Kriteria \& Bentuk Penilaian & Indikator Penilaian & Bobot Penilaian (\%) \\
(1) & (2) & (3) & (4) & (5) & (6) & (7) & (8) & (9) \\
1 & SCPMK0801-03601 & Kickoff proyek: pembentukan tim dan requirement gathering. & Modalitas: Blended Learning. Bentuk: Luring/Daring. Metode: Project Based Learning, sprint, kolaborasi tim, dan presentasi. & 1 x 3 x 50' tatap muka; 1 x 3 x 50' tugas/praktik mandiri. & Mahasiswa mengerjakan aktivitas terarah untuk kickoff proyek dan requirement gathering dan menunjukkan evidence hasil belajar. & Bentuk penilaian: Kickoff Proyek: Pembentukan Tim dan Requirement Gathering. Mengacu rubrik RTM. & Ketepatan implementasi; kualitas artefak proyek; kolaborasi dan dokumentasi. & 5 \\
2 & SCPMK0801-03601 & Dokumen perencanaan: WBS, jadwal, anggaran, risiko. & Modalitas: Blended Learning. Bentuk: Luring/Daring. Metode: Project Based Learning, sprint, kolaborasi tim, dan presentasi. & 1 x 3 x 50' tatap muka; 1 x 3 x 50' tugas/praktik mandiri. & Mahasiswa mengerjakan aktivitas terarah untuk penyusunan dokumen perencanaan proyek SI dan menunjukkan evidence hasil belajar. & Bentuk penilaian: Dokumen Perencanaan: WBS, Jadwal, Anggaran, Risiko. Mengacu rubrik RTM. & Ketepatan implementasi; kualitas artefak proyek; kolaborasi dan dokumentasi. & 10 \\
3 & SCPMK0802-03602, SCPMK1008-03605 & Arsitektur SI dan tech stack selection: mengintegrasikan hasil analisis kebutuhan ke desain arsitektur. & Modalitas: Blended Learning. Bentuk: Luring/Daring. Metode: Project Based Learning, sprint, kolaborasi tim, dan presentasi. & 1 x 3 x 50' tatap muka; 1 x 3 x 50' tugas/praktik mandiri. & Mahasiswa mengerjakan aktivitas terarah untuk arsitektur SI dan tech stack selection, termasuk mengintegrasikan hasil analisis dan desain ke tech stack, dan menunjukkan evidence hasil belajar. & Bentuk penilaian: Arsitektur SI dan Tech Stack Selection. Mengacu rubrik RTM. & Ketepatan implementasi; kualitas artefak proyek; kolaborasi dan dokumentasi; koherensi integrasi analisis-desain ke tech stack. & 2 \\
4 & SCPMK0802-03602 & Sprint planning dan setup environment. & Modalitas: Blended Learning. Bentuk: Luring/Daring. Metode: Project Based Learning, sprint, kolaborasi tim, dan presentasi. & 1 x 3 x 50' tatap muka; 1 x 3 x 50' tugas/praktik mandiri. & Mahasiswa mengerjakan aktivitas terarah untuk sprint planning dan setup environment dan menunjukkan evidence hasil belajar. & Bentuk penilaian: Sprint Planning dan Setup Environment. Mengacu rubrik RTM. & Ketepatan implementasi; kualitas artefak proyek; kolaborasi dan dokumentasi. & 3 \\
5 & SCPMK0802-03602, SCPMK0607-03604 & Sprint 1: implementasi fitur utama back-end. & Modalitas: Blended Learning. Bentuk: Luring/Daring. Metode: Project Based Learning, sprint, kolaborasi tim, dan presentasi. & 1 x 3 x 50' tatap muka; 1 x 3 x 50' tugas/praktik mandiri. & Mahasiswa mengerjakan aktivitas terarah untuk implementasi fitur utama back-end pada Sprint 1 dan menunjukkan evidence hasil belajar. & Bentuk penilaian: Sprint 1: Implementasi Fitur Utama Back-End. Mengacu rubrik RTM. & Ketepatan implementasi; kualitas artefak proyek; kolaborasi dan dokumentasi. & 5 \\
6 & SCPMK0607-03604 & Sprint 1: implementasi front-end dan integrasi. & Modalitas: Blended Learning. Bentuk: Luring/Daring. Metode: Project Based Learning, sprint, kolaborasi tim, dan presentasi. & 1 x 3 x 50' tatap muka; 1 x 3 x 50' tugas/praktik mandiri. & Mahasiswa mengerjakan aktivitas terarah untuk implementasi front-end dan integrasi pada Sprint 1 dan menunjukkan evidence hasil belajar. & Bentuk penilaian: Sprint 1: Implementasi Front-End dan Integrasi. Mengacu rubrik RTM. & Ketepatan implementasi; kualitas artefak proyek; kolaborasi dan dokumentasi. & 5 \\
7 & SCPMK0803-03603 & Sprint 1 review dan retrospective. & Modalitas: Blended Learning. Bentuk: Luring/Daring. Metode: Project Based Learning, sprint, kolaborasi tim, dan presentasi. & 1 x 3 x 50' tatap muka; 1 x 3 x 50' tugas/praktik mandiri. & Mahasiswa mengerjakan aktivitas terarah untuk sprint 1 review dan retrospective dan menunjukkan evidence hasil belajar. & Bentuk penilaian: Sprint 1 Review dan Retrospective. Mengacu rubrik RTM. & Ketepatan implementasi; kualitas artefak proyek; kolaborasi dan dokumentasi. & 5 \\
8 & SCPMK0803-03603 & UTS: milestone review dan demo Sprint 1-2. & Modalitas: Blended Learning. Bentuk: Luring/Daring. Metode: Project Based Learning, sprint, kolaborasi tim, dan presentasi. & 1 x 3 x 50' tatap muka; 1 x 3 x 50' tugas/praktik mandiri. & Mahasiswa mengerjakan aktivitas terarah untuk UTS milestone review dan demo Sprint 1-2 dan menunjukkan evidence hasil belajar. & Bentuk penilaian: UTS: Milestone Review dan Demo Sprint 1-2. Mengacu rubrik RTM. & Ketepatan implementasi; kualitas artefak proyek; kolaborasi dan dokumentasi. & 15 \\
9 & SCPMK0802-03602, SCPMK0607-03604 & Sprint 2: pengembangan fitur lanjutan. & Modalitas: Blended Learning. Bentuk: Luring/Daring. Metode: Project Based Learning, sprint, kolaborasi tim, dan presentasi. & 1 x 3 x 50' tatap muka; 1 x 3 x 50' tugas/praktik mandiri. & Mahasiswa mengerjakan aktivitas terarah untuk pengembangan fitur lanjutan pada Sprint 2 dan menunjukkan evidence hasil belajar. & Bentuk penilaian: Sprint 2: Pengembangan Fitur Lanjutan. Mengacu rubrik RTM. & Ketepatan implementasi; kualitas artefak proyek; kolaborasi dan dokumentasi. & 4 \\
10 & SCPMK0802-03602 & Sprint 2: testing dan debugging. & Modalitas: Blended Learning. Bentuk: Luring/Daring. Metode: Project Based Learning, sprint, kolaborasi tim, dan presentasi. & 1 x 3 x 50' tatap muka; 1 x 3 x 50' tugas/praktik mandiri. & Mahasiswa mengerjakan aktivitas terarah untuk testing dan debugging pada Sprint 2 dan menunjukkan evidence hasil belajar. & Bentuk penilaian: Sprint 2: Testing dan Debugging. Mengacu rubrik RTM. & Ketepatan implementasi; kualitas artefak proyek; kolaborasi dan dokumentasi. & 3 \\
11 & SCPMK0803-03603 & Sprint 2 review dan Sprint 3 planning. & Modalitas: Blended Learning. Bentuk: Luring/Daring. Metode: Project Based Learning, sprint, kolaborasi tim, dan presentasi. & 1 x 3 x 50' tatap muka; 1 x 3 x 50' tugas/praktik mandiri. & Mahasiswa mengerjakan aktivitas terarah untuk sprint 2 review dan sprint 3 planning dan menunjukkan evidence hasil belajar. & Bentuk penilaian: Sprint 2 Review dan Sprint 3 Planning. Mengacu rubrik RTM. & Ketepatan implementasi; kualitas artefak proyek; kolaborasi dan dokumentasi. & 3 \\
12 & SCPMK0607-03604 & Sprint 3: deployment preparation. & Modalitas: Blended Learning. Bentuk: Luring/Daring. Metode: Project Based Learning, sprint, kolaborasi tim, dan presentasi. & 1 x 3 x 50' tatap muka; 1 x 3 x 50' tugas/praktik mandiri. & Mahasiswa mengerjakan aktivitas terarah untuk deployment preparation pada Sprint 3 dan menunjukkan evidence hasil belajar. & Bentuk penilaian: Sprint 3: Deployment Preparation. Mengacu rubrik RTM. & Ketepatan implementasi; kualitas artefak proyek; kolaborasi dan dokumentasi. & 4 \\
13 & SCPMK0607-03604 & Sprint 3: UAT dan perbaikan. & Modalitas: Blended Learning. Bentuk: Luring/Daring. Metode: Project Based Learning, sprint, kolaborasi tim, dan presentasi. & 1 x 3 x 50' tatap muka; 1 x 3 x 50' tugas/praktik mandiri. & Mahasiswa mengerjakan aktivitas terarah untuk UAT dan perbaikan pada Sprint 3 dan menunjukkan evidence hasil belajar. & Bentuk penilaian: Sprint 3: UAT dan Perbaikan. Mengacu rubrik RTM. & Ketepatan implementasi; kualitas artefak proyek; kolaborasi dan dokumentasi. & 4 \\
14 & SCPMK0803-03603 & Sprint 3 review dan finalisasi. & Modalitas: Blended Learning. Bentuk: Luring/Daring. Metode: Project Based Learning, sprint, kolaborasi tim, dan presentasi. & 1 x 3 x 50' tatap muka; 1 x 3 x 50' tugas/praktik mandiri. & Mahasiswa mengerjakan aktivitas terarah untuk sprint 3 review dan finalisasi dan menunjukkan evidence hasil belajar. & Bentuk penilaian: Sprint 3 Review dan Finalisasi. Mengacu rubrik RTM. & Ketepatan implementasi; kualitas artefak proyek; kolaborasi dan dokumentasi. & 3 \\
15 & SCPMK0803-03603, SCPMK0607-03604, SCPMK0302-03606 & Penyusunan laporan akhir proyek dan konsolidasi hasil kerja tim. & Modalitas: Blended Learning. Bentuk: Luring/Daring. Metode: Project Based Learning, sprint, kolaborasi tim, dan presentasi. & 1 x 3 x 50' tatap muka; 1 x 3 x 50' tugas/praktik mandiri. & Mahasiswa mengerjakan aktivitas terarah untuk penyusunan laporan akhir proyek, termasuk mengelola dinamika tim dan menyusun penyampaian hasil kerja tim, dan menunjukkan evidence hasil belajar. & Bentuk penilaian: Penyusunan Laporan Akhir Proyek. Mengacu rubrik RTM. & Ketepatan implementasi; kualitas artefak proyek; kolaborasi dan dokumentasi; koordinasi tim dalam penyusunan laporan. & 10 \\
16 & SCPMK0607-03604 & Persiapan demo dan presentasi. & Modalitas: Blended Learning. Bentuk: Luring/Daring. Metode: Project Based Learning, sprint, kolaborasi tim, dan presentasi. & 1 x 3 x 50' tatap muka; 1 x 3 x 50' tugas/praktik mandiri. & Mahasiswa mengerjakan aktivitas terarah untuk persiapan demo dan presentasi dan menunjukkan evidence hasil belajar. & Bentuk penilaian: Persiapan Demo dan Presentasi. Mengacu rubrik RTM. & Ketepatan implementasi; kualitas artefak proyek; kolaborasi dan dokumentasi. & 5 \\
17 & SCPMK0607-03604, SCPMK0803-03603, SCPMK0302-03606 & UAS: demo produk dan defense proyek SI, termasuk penyampaian hasil kerja tim kepada pemangku kepentingan. & Modalitas: Blended Learning. Bentuk: Luring/Daring. Metode: Project Based Learning, sprint, kolaborasi tim, dan presentasi. & 1 x 3 x 50' tatap muka; 1 x 3 x 50' tugas/praktik mandiri. & Mahasiswa mengerjakan aktivitas terarah untuk UAS demo produk dan defense proyek SI, termasuk penyampaian hasil kerja tim kepada pemangku kepentingan, dan menunjukkan evidence hasil belajar. & Bentuk penilaian: UAS: Demo Produk dan Defense Proyek SI. Mengacu rubrik RTM. & Ketepatan implementasi; kualitas artefak proyek; kolaborasi dan dokumentasi; profesionalisme penyampaian hasil kepada pemangku kepentingan. & 15 \\
\end{longtblr}
